\documentclass[10pt,a4paper]{amsart}
\usepackage{amsmath,amssymb,amsthm,graphicx}
% Retained minimal preamble for the Stacks Project chapter more-morphisms.tex (More on Morphisms).
% Verbatim excerpt of the chapter preamble: Xy-pic setup (native preamble lines 14-18) and the native
% theorem declarations with the mathematical macros (native preamble lines 154-217). Copied mechanically;
% nothing added, removed or redefined. The wrapper class and the amsmath/amssymb/amsthm/graphicx base
% packages are supplied by the assembler, so the class/externaldocument part of the native preamble is omitted.
% For commutative diagrams we use Xy-pic
\usepackage[all]{xy}

% We use 2cell for 2-commutative diagrams.
\xyoption{2cell}
% Theorem environments.
%
\theoremstyle{plain}
\newtheorem{theorem}[subsection]{Theorem}
\newtheorem{proposition}[subsection]{Proposition}
\newtheorem{lemma}[subsection]{Lemma}

\theoremstyle{definition}
\newtheorem{definition}[subsection]{Definition}
\newtheorem{example}[subsection]{Example}
\newtheorem{exercise}[subsection]{Exercise}
\newtheorem{situation}[subsection]{Situation}

\theoremstyle{remark}
\newtheorem{remark}[subsection]{Remark}
\newtheorem{remarks}[subsection]{Remarks}

\numberwithin{equation}{subsection}

% Macros
%
\def\lim{\mathop{\mathrm{lim}}\nolimits}
\def\colim{\mathop{\mathrm{colim}}\nolimits}
\def\Spec{\mathop{\mathrm{Spec}}}
\def\Hom{\mathop{\mathrm{Hom}}\nolimits}
\def\Ext{\mathop{\mathrm{Ext}}\nolimits}
\def\SheafHom{\mathop{\mathcal{H}\!\mathit{om}}\nolimits}
\def\SheafExt{\mathop{\mathcal{E}\!\mathit{xt}}\nolimits}
\def\Sch{\mathit{Sch}}
\def\Mor{\mathop{\mathrm{Mor}}\nolimits}
\def\Ob{\mathop{\mathrm{Ob}}\nolimits}
\def\Sh{\mathop{\mathit{Sh}}\nolimits}
\def\NL{\mathop{N\!L}\nolimits}
\def\CH{\mathop{\mathrm{CH}}\nolimits}
\def\proetale{{pro\text{-}\acute{e}tale}}
\def\etale{{\acute{e}tale}}
\def\QCoh{\mathit{QCoh}}
\def\Ker{\mathop{\mathrm{Ker}}}
\def\Im{\mathop{\mathrm{Im}}}
\def\Coker{\mathop{\mathrm{Coker}}}
\def\Coim{\mathop{\mathrm{Coim}}}

% Boxtimes
%
\DeclareMathSymbol{\boxtimes}{\mathbin}{AMSa}{"02}

%
% Macros for moduli stacks/spaces
%
\def\QCohstack{\mathcal{QC}\!\mathit{oh}}
\def\Cohstack{\mathcal{C}\!\mathit{oh}}
\def\Spacesstack{\mathcal{S}\!\mathit{paces}}
\def\Quotfunctor{\mathrm{Quot}}
\def\Hilbfunctor{\mathrm{Hilb}}
\def\Curvesstack{\mathcal{C}\!\mathit{urves}}
\def\Polarizedstack{\mathcal{P}\!\mathit{olarized}}
\def\Complexesstack{\mathcal{C}\!\mathit{omplexes}}
% \Pic is the operator that assigns to X its picard group, usage \Pic(X)
% \Picardstack_{X/B} denotes the Picard stack of X over B
% \Picardfunctor_{X/B} denotes the Picard functor of X over B
\def\Pic{\mathop{\mathrm{Pic}}\nolimits}
\def\Picardstack{\mathcal{P}\!\mathit{ic}}
\def\Picardfunctor{\mathrm{Pic}}
\def\Deformationcategory{\mathcal{D}\!\mathit{ef}}
\title{Universal openness of Spec(B) over Spec(A) tested on one explicit principal open for a finite generating set}
\author{The Stacks Project Authors}
\date{stacks-project:more-morphisms.tex:local-snapshot-20261011; GNU FDL 1.2 or later}
\begin{document}
\maketitle
\noindent\textit{Editorial note.} Statement span is the whole lemma environment lines 21846-21852; proof span is lines 21854-21875, the immediately following proof environment.\par\medskip
\noindent\textit{Editorial note.} The proof contains two xymatrix displays; the retained preamble provides the Xy-pic setup, so no extra\_packages are declared.\par\medskip
\noindent\textit{Editorial note.} No external\_references are needed: the proof cites no label and no bibliography entry.\par\medskip
\noindent\textit{Editorial note.} The proof text is retained verbatim, including the mild imprecision in the original (the statement writes x\_1,...,x\_d while the proof writes x\_1,...,x\_n); no correction was applied.\par\medskip
\noindent\textit{Editorial note.} extra\_packages is empty on purpose: the retained preamble already loads xy with options and a second declaration clashes.\par\medskip
\noindent\textit{Allowed standard prerequisites (editorial).}\begin{itemize}
\item Explicit prerequisite boundary. Base change of morphisms of schemes, the fact that principal opens D(g) form a basis of the topology, polynomial algebras and the cartesian square Spec(B[x\_1..x\_n]) / Spec(A[x\_1..x\_n]) against Spec(B') / Spec(A'). All are standard; the reduction itself is retained.
\end{itemize}\medskip
\begin{lemma}
\label{lemma-characterize-universally-open-finite}
Let $A \to B$ be a ring map. Say $B$ is generated as an $A$-module by
$b_1, \ldots, b_d \in B$. Set $h = \sum x_ib_i \in B[x_1, \ldots, x_d]$.
Then $\Spec(B) \to \Spec(A)$ is universally open if and only if the image of
$D(h)$ in $\Spec(A[x_1, \ldots, x_d])$ is open.
\end{lemma}

\begin{proof}
If $\Spec(B) \to \Spec(A)$ is universally open, then of course
the image of $D(h)$ is open. Conversely, assume the image $U$ of
$D(h)$ is open. Let $A \to A'$ be a ring map. It suffices to show
that the image of any principal open $D(g) \subset \Spec(A' \otimes_A B)$
in $\Spec(A')$ is open. We may write
$g = \sum_{i = 1, \ldots, d} a'_i \otimes b_i$ for some $a'_i \in A'$.
Let $A[x_1, \ldots, x_n] \to A'$ be the $A$-algebra homomorphism
sending $x_i$ to $a'_i$. This also induces a $B$-algebra
homomorphism $B[x_1, \ldots, x_n] \to A' \otimes_A B$ sending $h$ to $g$.
Since
$$
\xymatrix{
\Spec(B[x_1, \ldots, x_n]) \ar[d] &
\Spec(B') \ar[l] \ar[d] \\
\Spec(A[x_1, \ldots, x_n]) &
\Spec(A') \ar[l]
}
$$
is cartesian the image of $D(g)$ in $\Spec(A')$ is equal to the
inverse image of $U$ in $\Spec(A')$ and hence open.
\end{proof}

\end{document}
